% !TEX root = main.tex

% --- Présentation de l'ENI-ABT ---
\chapter*{Présentation de l'ENI-ABT}
\addcontentsline{toc}{chapter}{Présentation de l'ENI-ABT}

\begin{tcolorbox}[colback=white, colframe=black, sharp corners, halign=center, boxrule=1.5pt]
    \textbf{\Large\textcolor{macouleur}{Présentation de L'ENI-ABT}}
\end{tcolorbox}

\vspace{0.5cm}
Créée le 14 avril 1939 l'\textbf{Ecole Nationale d'Ingénieurs – Abderhamane Baba Touré} (ENI-ABT) a subi plusieurs mutations allant de l'\textbf{Ecole des Travaux Publics} formant des dessinateurs jusqu'à la formation des ingénieurs de conception dans plusieurs domaines. Elle commença comme l'\textbf{Ecole des Travaux Publics} (ETP) de l'A.O.P destinée à former des adjoints techniques en Travaux Publics. La première réforme de l'enseignement érigea l'ETP en \textbf{Ecole Nationale d'Ingénieurs} (ENI) en 1962, en Septembre 1973 le recrutement à l'ENI fut conditionné au Baccalauréat, pour la formation des ingénieurs des sciences appliquées en industrie, génie civil, géologie et en topographie. En 2002 l'ENI, à l'instar des autres grandes écoles (IPR/IFRA, ENSUP), a été directement rattachée à la Direction Nationale de l'Enseignement Supérieur et de la Recherche Scientifique (DNESRS) suivant l'Ordonnance N° 02-054/P-RM du 04 juin 2002. Elle a alors été régie dans son organisation et son fonctionnement par le décret N° 96-378/P-RM du 31 décembre 1996. En juin 2006 l'ENI fut baptisée au nom de son premier Directeur Malien Abderhamane Baba Touré (ENI\textendash ABT).

\vspace{0.5cm}
L'ENI-ABT occupe une place centrale dans la politique industrielle du Mali. En formant des ingénieurs capables de concevoir, de fabriquer et de maintenir des systèmes mécaniques complexes, elle répond directement aux besoins d'une économie en pleine mutation vers l'industrialisation. Le présent Projet de Fin d'Études (PFE) s'inscrit dans cette vision : démontrer qu'il est possible de concevoir et de réaliser localement, à l'ENI-ABT, une machine-outil de précision à 5 axes, en exploitant des technologies de l'Industrie 4.0, sans dépendance aux systèmes propriétaires étrangers. L'ENI-ABT dispose à ce jour d'un atelier mécanique équipé de tours parallèles, de fraiseuses conventionnelles et d'une imprimante 3D FDM, ce qui constitue un socle technique sur lequel ce projet vient greffer une dimension CNC numérique avancée.

\vspace{1cm}
\begin{figure}[htbp]
    \centering
    \caption{Vues des bâtiments de l'ENI-ABT, Bamako, Mali}
    \label{fig:eni_abt}
\end{figure}

\newpage
% --- Structuration pédagogique ---
\section*{\textcolor{macouleur}{Structuration pédagogique}}

Depuis la rentrée 2014-2015, l'ENI-ABT applique le système LMD à travers une nouvelle structuration comportant six DER pour six filières de départ :

\vspace{0.5cm}
\noindent
\begin{minipage}[t]{0.48\textwidth}
    \begin{tcolorbox}[colback=white, colframe=gray, boxrule=1pt, sharp corners, halign=center]
        \textbf{L'ENI-ABT avec le système}\\
        \textbf{«CLASSIQUE»}
    \end{tcolorbox}
    \vspace{0.3cm}
    \textbf{Deux niveaux de formations :}
    \begin{itemize}[label=\raisebox{0.25ex}{\tiny$\blacktriangleright$}, itemsep=2pt]
        \item Cycle Ingénieur ;
        \item Cycle Technicien Supérieur.
    \end{itemize}
    \vspace{0.5cm}
    \textbf{Quatre DER :}
    \begin{itemize}[label=\raisebox{0.25ex}{\tiny$\blacktriangleright$}, itemsep=2pt]
        \item Génie civil ;
        \item Génie Industriel ;
        \item Géodésie ;
        \item Géologie.
    \end{itemize}
\end{minipage}
\hfill
\vrule
\hfill
\begin{minipage}[t]{0.48\textwidth}
    \begin{tcolorbox}[colback=gray!10, colframe=gray, boxrule=1pt, sharp corners, halign=center]
        \textbf{L'ENI-ABT avec le système}\\
        \textbf{«LMD»}
    \end{tcolorbox}
    \vspace{0.3cm}
    \textbf{Trois niveaux de formations :}
    \begin{itemize}[label=\raisebox{0.25ex}{\tiny$\blacktriangleright$}, itemsep=2pt]
        \item Licence ;
        \item Master ;
        \item Doctorat.
    \end{itemize}
    \vspace{0.5cm}
    \textbf{Six DER :}
    \begin{itemize}[label=\raisebox{0.25ex}{\tiny$\blacktriangleright$}, itemsep=2pt]
        \item Génie civil ;
        \item Génie Electrique/Industrie ;
        \item Génie Mécanique et Energie ;
        \item Génie Informatique et Télécommunication ;
        \item Géodésie/Sciences Géographiques.
    \end{itemize}
\end{minipage}

\vspace{1cm}
\noindent L'ENI-ABT mène également des activités :
\begin{itemize}[label=\raisebox{0.25ex}{\tiny$\blacktriangleright$}]
    \item De formation continue diplômante et qualifiante à travers son Service de la Formation Continue (SFC) ;
    \item D'expertise et de production à travers son Unité d'Expertise et de Production (UEP).
\end{itemize}

\newpage
% ===================================================
% --- Sigles et Abréviations ---
\chapter*{Sigles et Abréviations}
\addcontentsline{toc}{chapter}{Sigles et Abréviations}

\renewcommand{\arraystretch}{1.25}
\begin{longtable}{|l|p{10.5cm}|}
\hline
\rowcolor{macouleur!20}
\textbf{Sigle / Abréviation} & \textbf{Signification} \\
\hline
\endfirsthead
\hline
\rowcolor{macouleur!20}
\textbf{Sigle / Abréviation} & \textbf{Signification} \\
\hline
\endhead
\hline
\endfoot
\textbf{AMDEC} & Analyse des Modes de Défaillance, de leurs Effets et de leur Criticité \\
\hline
\textbf{AOP} & Afrique Occidentale Protectorat \\
\hline
\textbf{API} & Application Programming Interface \\
\hline
\textbf{BOM} & Bill of Materials (Nomenclature financière de projet) \\
\hline
\textbf{CAO} & Conception Assistée par Ordinateur \\
\hline
\textbf{CEM} & Compatibilité ÉlectroMagnétique (EMC en anglais) \\
\hline
\textbf{CNC} & Computer Numerical Control (Commande Numérique par Calculateur) \\
\hline
\textbf{D-H} & Denavit-Hartenberg (Paramètres conventionnels de cinématique des robots) \\
\hline
\textbf{DDA} & Digital Differential Analyzer (Algorithme de génération de trajectoire) \\
\hline
\textbf{DER} & Département d'Enseignement et de Recherche \\
\hline
\textbf{DPU} & Data Processing Unit (Unité de Traitement des Données) \\
\hline
\textbf{EEPROM} & Electrically Erasable Programmable Read-Only Memory \\
\hline
\textbf{EMI} & ElectroMagnetic Interference (Perturbation électromagnétique) \\
\hline
\textbf{ENI-ABT} & École Nationale d'Ingénieurs – Abderhamane Baba Touré \\
\hline
\textbf{E-STOP} & Emergency Stop (Arrêt d'Urgence) \\
\hline
\textbf{FAO} & Fabrication Assistée par Ordinateur \\
\hline
\textbf{FAST} & Function Analysis System Technique (Diagramme d'analyse fonctionnelle) \\
\hline
\textbf{FCFA} & Franc de la Communauté Financière Africaine \\
\hline
\textbf{FEA} & Finite Element Analysis (Analyse par Éléments Finis) \\
\hline
\textbf{FFF} & Fused Filament Fabrication (impression 3D par dépôt de matière) \\
\hline
\textbf{FreeRTOS} & Free Real-Time Operating System (OS temps réel open-source) \\
\hline
\textbf{FSM} & Finite State Machine (Machine à États Finis) \\
\hline
\textbf{GPIO} & General Purpose Input/Output (Broches d'entrée/sortie) \\
\hline
\textbf{GRBL} & G-code Real-time Bootstrap Loader (firmware CNC open-source) \\
\hline
\textbf{HB} & Dureté Brinell \\
\hline
\textbf{IDE} & Integrated Development Environment (Environnement de développement) \\
\hline
\textbf{IHM} & Interface Homme-Machine \\
\hline
\textbf{IoT} & Internet of Things (Internet des Objets) \\
\hline
\textbf{ISR} & Interrupt Service Routine (Routine d'interruption matérielle) \\
\hline
\textbf{LDO} & Low Drop-Out (Régulateur de tension linéaire) \\
\hline
\textbf{LUT} & Look-Up Table (Table de correspondance pré-calculée) \\
\hline
\textbf{MISRA} & Motor Industry Software Reliability Association \\
\hline
\textbf{NEMA} & National Electrical Manufacturers Association \\
\hline
\textbf{OEE} & Overall Equipment Effectiveness (Taux de Rendement Synthétique) \\
\hline
\textbf{PCB} & Printed Circuit Board (Circuit Imprimé) \\
\hline
\textbf{PFE} & Projet de Fin d'Études \\
\hline
\textbf{PME} & Petites et Moyennes Entreprises \\
\hline
\textbf{PSU} & Power Supply Unit (Alimentation à Découpage) \\
\hline
\textbf{PWM} & Pulse Width Modulation (Modulation de Largeur d'Impulsion) \\
\hline
\textbf{RDM} & Résistance des Matériaux \\
\hline
\textbf{ROI} & Return on Investment (Retour sur Investissement) \\
\hline
\textbf{RTOS} & Real-Time Operating System \\
\hline
\textbf{SMPS} & Switch-Mode Power Supply (Alimentation à découpage industrielle) \\
\hline
\textbf{SoC} & System on Chip (Système sur puce intégrée) \\
\hline
\textbf{TCP} & Tool Center Point (Point de contact de l'outil) \\
\hline
\textbf{TTL} & Transistor-Transistor Logic (Logique $5\text{V}$) \\
\hline
\textbf{UDP} & User Datagram Protocol (Protocole réseau sans connexion) \\
\hline
\textbf{UML} & Unified Modeling Language \\
\hline
\textbf{USB} & Universal Serial Bus \\
\hline
\textbf{VFD} & Variable Frequency Drive (Variateur de fréquence pour broche) \\
\hline
\textbf{WBS} & Work Breakdown Structure (Structure de découpage du travail) \\
\hline
\end{longtable}
\newpage
% ===================================================
% --- NOTATIONS ET SYMBOLES MATHÉMATIQUES ---
\chapter*{Notations et Symboles Mathématiques}
\addcontentsline{toc}{chapter}{Notations et Symboles Mathématiques}

Le tableau suivant recense les principaux symboles mathématiques utilisés dans les chapitres de calcul, accompagnés de leur signification physique et de leur unité dans le Système International (SI).

\vspace{0.3cm}
\renewcommand{\arraystretch}{1.3}
\begin{center}
\begin{tabular}{|c|p{8cm}|c|}
\hline
\rowcolor{macouleur!20}
\textbf{Symbole} & \textbf{Signification} & \textbf{Unité} \\
\hline
$F_c$ & Force de coupe tangentielle & N \\
\hline
$F_f$ & Force d'avance (axiale) & N \\
\hline
$F_p$ & Force de pénétration (radiale) & N \\
\hline
$F_{ax}$ & Force axiale sur la vis à billes & N \\
\hline
$\vec{R}$ & Résultante des forces de coupe (torseur) & N \\
\hline
$M_f$ & Moment de flexion & N·mm \\
\hline
$M_{renversement}$ & Moment de renversement du chariot Y (Table Trunnion) & N·m \\
\hline
$\sigma_{eq}$ & Contrainte équivalente de Von Mises & MPa \\
\hline
$\sigma_{max}$ & Contrainte normale maximale en flexion & MPa \\
\hline
$\tau$ & Contrainte de cisaillement & MPa \\
\hline
$R_e$ & Limite d'élasticité du matériau & MPa \\
\hline
$R_m$ & Résistance maximale à la rupture & MPa \\
\hline
$E$ & Module de Young (rigidité) & MPa \\
\hline
$I_{yy}$ & Moment quadratique de la section transversale & mm$^4$ \\
\hline
$f_{max}$ & Flèche maximale (déflexion statique) & mm \\
\hline
$F_k$ & Charge critique de flambage (critère d'Euler) & N \\
\hline
$k_L$ & Facteur d'appui selon le montage de la vis & --- \\
\hline
$s$ & Coefficient de sécurité statique & --- \\
\hline
$C_r$ & Couple résistant continu de la vis à billes & N·m \\
\hline
$C_J$ & Couple d'inertie lors des accélérations & N·m \\
\hline
$P_h$ & Pas de la vis à billes & mm/tour \\
\hline
$\eta$ & Rendement de la transmission par vis à billes & \% \\
\hline
$\Delta\theta$ & Erreur angulaire du moteur (Backlash) & ° ou rad \\
\hline
$\Delta L$ & Erreur linéaire tangentielle (effet levier Trunnion) & mm \\
\hline
$R$ & Rayon du plateau Trunnion & mm \\
\hline
$H$ & Hauteur du berceau Trunnion & mm \\
\hline
$L_y, L_z$ & Offsets géométriques de Denavit-Hartenberg & mm \\
\hline
$\theta_A, \theta_B$ & Angles des articulations rotatives Axe A et Axe B & ° \\
\hline
$\,^0T_n$ & Matrice de transformation homogène globale (produit D-H) & --- \\
\hline
$f_{step}$ & Fréquence des impulsions STEP vers les drivers & Hz \\
\hline
$k_c$ & Pression spécifique de coupe du matériau & N/mm$^2$ \\
\hline
$V_c$ & Vitesse de coupe outil & m/min \\
\hline
$f_z$ & Avance par dent de la fraise & mm/dent \\
\hline
$a_p$ & Profondeur de passe axiale & mm \\
\hline
$a_e$ & Profondeur de passe radiale (engagement) & mm \\
\hline
$D_c$ & Diamètre de l'outil coupant (fraise) & mm \\
\hline
$P_{PSU}$ & Puissance nominale de l'alimentation à découpage & W \\
\hline
$C_{bulk}$ & Condensateur de découplage (anti Brown-out) & F (µF) \\
\hline
$L_{10}$ & Durée de vie nominale des roulements (ISO 281) & heures (h) \\
\hline
$C_{dyn}$ & Charge dynamique de base des roulements & kN \\
\hline
$P_{equiv}$ & Charge équivalente sur le roulement & N \\
\hline
$Ra$ & Rugosité moyenne arithmétique de surface & µm \\
\hline
$C_{pk}$ & Indice de capabilité machine (ISO 22514) & --- \\
\hline
\end{tabular}
\end{center}
